%!TeX root=MemoriaTFG.tex

\chapter{Conjunto de datos y ontología}
\label{cap:datos}

Este capítulo describe el conjunto de datos sobre el que se realiza
el trabajo experimental y, de forma resumida, el corpus propio
resultante del trabajo. El núcleo experimental son los doce datasets
externos públicos descritos por familia, sobre los que se evalúan
los modelos del estado del arte. El corpus propio DermapixelAI 1.0,
derivado del archivo Dermapixel y desarrollado en paralelo al
presente trabajo en el marco de la modernización tecnológica del
archivo descrito en el capítulo~\ref{cap:intro}, se introduce en la
misma sección que los datasets externos y se documenta de forma
exhaustiva en el anexo~\ref{anexo:pipeline}; no es un dataset
preexistente sobre el que se evalúan modelos, sino una contribución
del propio proyecto a la comunidad investigadora hispanohablante.

Tras la presentación de los datasets se introduce la ontología
jerárquica L1/L2/L3 utilizada para armonizar las taxonomías
diagnósticas nativas heterogéneas. El capítulo cierra con un
inventario de las limitaciones del conjunto de datos sobre el que
opera la evaluación experimental del trabajo.

\section{Vista general del material de trabajo}
\label{sec:vista-general-datos}

El TFG opera sobre material procedente de tres orígenes distintos:
datasets públicos publicados por la comunidad científica, datasets
institucionales liberados parcialmente para investigación, y el
corpus propio del autor construido sobre el archivo Dermapixel. La
tabla~\ref{tab:datasets-resumen} resume las características
principales de los doce datasets externos sobre los que se realiza
la evaluación experimental del trabajo. El corpus propio
DermapixelAI 1.0 se introduce de forma resumida en la
sección~\ref{sec:dermapixel-summary}, junto a los demás conjuntos
descritos por familia, y su pipeline de construcción y
caracterización cuantitativa completa se desarrollan en el
anexo~\ref{anexo:pipeline}, dado que el corpus no constituye objeto
de evaluación del trabajo sino resultado del mismo.

\begin{table}[H]
\centering
\footnotesize
\setlength{\tabcolsep}{4pt}
\renewcommand{\arraystretch}{0.95}
\caption{Resumen del material de trabajo. ``$N_{\mathrm{test}}$''
denota el tamaño nominal del split de evaluación canónico empleado en
linear probing; ``Clases'' indica el número de categorías
diagnósticas o conceptos según la convención propia del dataset.
Un asterisco (\textsuperscript{$\ast$}) en ``Uso principal'' marca
los datasets con caveat de leakage frente al corpus de
preentrenamiento Derm1M.}
\label{tab:datasets-resumen}
\begin{tabular}{lrcll}
\hline
Dataset & $N_{\mathrm{test}}$ & Clases & Modalidad & Uso principal \\
\hline
HAM10000~\cite{ham10000}             & 1\,232  & 7          & Dermoscopia      & LP, FT, ZS, LLMs \\
BCN20000~\cite{bcn20000}             & 1\,242  & 9          & Dermoscopia      & LP \\
PAD-UFES-20~\cite{padufes}           & 461     & 6          & Clínica (móvil)  & LP, FT \\
DDI~\cite{ddi}                       & 137     & 2          & Clínica          & LP, FT, equidad \\
Dermnet~\cite{dermnet}               & 4\,002  & 23         & Clínica (atlas)  & LP\textsuperscript{$\ast$} \\
Derm7pt clínico~\cite{kawahara2019}  & 168     & 2          & Clínica          & LP \\
Derm7pt dermo~\cite{kawahara2019}    & 225     & 2          & Dermoscopia      & LP \\
HIBA                                 & 334     & 2          & Dermoscopia      & LP\textsuperscript{$\ast$} \\
MSKCC                                & 1\,664  & 2          & Dermoscopia      & LP, FT\textsuperscript{$\ast$} \\
WSI patches                          & 12\,354 & 16         & Histopatología   & LP \\
ISIC2018~\cite{isic2018}             & 2\,594  & binaria    & Dermoscopia      & Segmentación (LoRA) \\
Fitzpatrick17k~\cite{fitzpatrick17k} & 16\,577 & 114 + 6 FP & Clínica          & Equidad fototipo \\
SkinCon~\cite{skincon}               & 3\,230  & 48 conc.   & Clínica          & CBM sobre SAE \\
\hline
\end{tabular}
\end{table}

La columna ``$N_{\mathrm{test}}$'' reúne las cifras nominales del
split de evaluación que el presente trabajo emplea, no necesariamente
el tamaño total del dataset. La columna ``Uso principal'' anticipa el
protocolo experimental asociado a cada dataset, cuya descripción
formal se desarrolla en el capítulo~\ref{cap:experimental}. Tres
datasets aparecen marcados con un \emph{caveat de leakage} frente al
corpus de preentrenamiento Derm1M (Dermnet, HIBA, MSKCC). En este
contexto, el término \emph{caveat} designa una advertencia
metodológica que debe acompañar al uso del dataset; \emph{leakage}
designa el solapamiento entre las imágenes del split de evaluación y
las imágenes con las que el modelo se preentrenó. La consecuencia
operativa es que las métricas obtenidas sobre un dataset afectado por
leakage no reflejan la capacidad del modelo de generalizar a material
nuevo, sino que pueden recoger, total o parcialmente, su capacidad de
reconocer material previamente visto durante el preentrenamiento. El
tratamiento detallado de este caveat se desarrolla en el
capítulo~\ref{cap:resultados}.

\section{Datasets externos por familia}
\label{sec:datasets-externos}

Los datasets externos se agrupan por familia de modalidad y
finalidad. Las cifras presentadas en cada subsección corresponden a
los valores publicados por los responsables de cada conjunto y se
emplean tal como están en los splits canónicos.

\subsection{Datasets dermatoscópicos}
\label{sec:datasets-dermo}

\textbf{HAM10000}~\cite{ham10000} (\emph{Human Against Machine with
10\,000 training images}) es el dataset dermatoscópico de referencia
para clasificación multiclase desde su publicación. Reúne 10\,015
imágenes dermatoscópicas de siete clases (\emph{nv}, \emph{mel},
\emph{bkl}, \emph{bcc}, \emph{akiec}, \emph{vasc}, \emph{df}) con un
fuerte desbalance hacia \emph{nv}. El split oficial utilizado en este
trabajo dispone de 1\,232 imágenes en test.

\textbf{BCN20000}~\cite{bcn20000} (\emph{Barcelona Dataset for Skin
Lesion Classification}) extiende el espectro diagnóstico a nueve
clases sobre dermoscopias del Hospital Clínic de Barcelona. El test
empleado en este trabajo contiene 1\,242 imágenes. Su distribución de
clases incluye, además de las siete categorías de HAM10000, dos
adicionales asociadas a lesiones no clasificables y a entidades poco
representadas en HAM.

\textbf{Derm7pt} (variante dermatoscópica)~\cite{kawahara2019} forma
parte del corpus \emph{Seven-Point Checklist} y constituye, junto a
su variante clínica, uno de los pocos datasets con anotaciones
multitarea ligadas a criterios diagnósticos estructurados. La
variante dermatoscópica del split de evaluación dispone de 225
imágenes en una formulación binaria melanoma/no-melanoma.

\textbf{HIBA} es un dataset dermatoscópico originario del Hospital
Italiano de Buenos Aires, liberado parcialmente para la comunidad
investigadora a través del ISIC Archive. El test utilizado contiene
334 imágenes en formulación binaria. Aparece marcado con caveat de
leakage frente a Derm1M.

\textbf{MSKCC} es el subconjunto dermatoscópico del Memorial Sloan
Kettering Cancer Center, distribuido a través del ISIC Archive. El
test empleado dispone de 1\,664 imágenes en formulación binaria
melanoma/benigno. También aparece con caveat de leakage frente a
Derm1M.

\textbf{ISIC2018}~\cite{isic2018}, en su tarea de segmentación
binaria de lesión, proporciona 2\,594 imágenes dermatoscópicas con
máscara binaria por imagen. Es el dataset canónico de referencia para
fine-tuning de modelos de segmentación en dermatología; en el
presente trabajo se emplea para el fine-tuning con LoRA de SAM2.1
descrito en el capítulo~\ref{cap:experimental}.

\subsection{Datasets de fotografía clínica}
\label{sec:datasets-clinicos}

\textbf{PAD-UFES-20}~\cite{padufes} (\emph{Patient and Dataset of
Skin Lesions, UFES}) recoge fotografías clínicas tomadas con
teléfono móvil en consultas comunitarias del estado brasileño de
Espírito Santo, cubriendo seis clases diagnósticas. Se trata del
material que mejor reproduce las condiciones reales de imagen
clínica generada por atención primaria: variabilidad de iluminación,
encuadre y dispositivo de captura. El test utilizado dispone de 461
imágenes.

\textbf{DDI}~\cite{ddi} (\emph{Diverse Dermatology Images}) es un
dataset curado específicamente para evaluar equidad: contiene 656
imágenes clínicas con representación equilibrada de fototipos I-II,
III-IV y V-VI, anotadas con una formulación binaria maligna/benigna.
El test utilizado contiene 137 imágenes. Su tamaño reducido y su
diversidad lo convierten en una referencia para análisis de
robustez por subgrupo cutáneo.

\textbf{Dermnet}~\cite{dermnet} es un dataset clínico construido a
partir del atlas dermatológico Dermnet NZ, con 23 categorías
diagnósticas y 4\,002 imágenes en test. Aparece con caveat de
leakage frente a Derm1M: parte de su material se solapa con el corpus
de preentrenamiento.

\textbf{Derm7pt} (variante clínica)~\cite{kawahara2019} es la
contraparte clínica del corpus \emph{Seven-Point Checklist}. El test
contiene 168 imágenes en formulación binaria melanoma/no-melanoma.
La existencia de las dos variantes (clínica y dermatoscópica) sobre
un mismo conjunto de lesiones permite contrastes intra-paciente
entre modalidades.

\subsection{Datasets histopatológicos}
\label{sec:datasets-histo}

\textbf{WSI patches} es un conjunto de \emph{parches} (recortes
rectangulares de tamaño fijo) extraídos de \emph{Whole-Slide Images}
de piezas histopatológicas dermatológicas, distribuido como benchmark
de evaluación para encoders fundacionales. El test utilizado en el
presente trabajo contiene 12\,354 parches sobre 16 clases
histopatológicas, derivadas de piezas etiquetadas por patólogos.

\subsection{Datasets especializados}
\label{sec:datasets-especiales}

Dos datasets se incorporan al trabajo no como benchmarks generales de
clasificación sino para evaluar dimensiones específicas del
rendimiento.

\textbf{Fitzpatrick17k}~\cite{fitzpatrick17k} agrupa 16\,577 imágenes
clínicas anotadas conjuntamente con diagnóstico (114 patologías
distribuidas en tres grandes familias etiológicas) y fototipo cutáneo
Fitzpatrick (escala de seis niveles, I-VI). Su finalidad en el
presente trabajo es la evaluación de equidad por fototipo descrita en
el capítulo~\ref{cap:experimental}.

\textbf{SkinCon}~\cite{skincon} es un dataset clínico de 3\,230
imágenes anotadas densamente con 48 conceptos clínicos visuales por
varios dermatólogos independientes. Su finalidad en el presente
trabajo es servir como ground-truth conceptual para la evaluación de
los Concept Bottleneck Models construidos sobre las representaciones
del Sparse Autoencoder.

\subsection{Corpus propio del trabajo: DermapixelAI 1.0}
\label{sec:dermapixel-summary}

\textbf{DermapixelAI 1.0} es el corpus propio resultante del
trabajo, construido a partir del archivo del blog Dermapixel
mantenido por la Dra.~R.~Taberner y desarrollado en paralelo al
TFG en el marco de la modernización tecnológica del archivo
descrita en el capítulo~\ref{cap:intro}. La versión 1.0 contiene
1\,089 imágenes vinculadas a 672 casos clínicos con texto narrativo
asociado en castellano (mediana de 223 palabras por caso),
mayoritariamente fotografía clínica (1\,036) con un subconjunto
dermatoscópico reducido (49) y unas pocas imágenes de otras
modalidades. Las imágenes están etiquetadas conforme a la ontología
jerárquica L1/L2/L3 descrita en la sección~\ref{sec:ontologia} y
particionadas en train/val/test bajo regla \emph{case-aware} (891
/ 157 / 41), con 97{,}93\,\% del corpus validado mediante mapeo
ontológico experto.

DermapixelAI 1.0 no se concibe como benchmark de evaluación del
presente TFG sino como contribución del proyecto a la comunidad
investigadora y, por extensión, al sistema sanitario. Su finalidad
principal es servir de base para investigaciones posteriores sobre
dermatología clínica en castellano ---en particular las orientadas
a la integración del texto clínico como componente del aprendizaje
multimodal--- y habilitar la reutilización del archivo Dermapixel
en proyectos académicos y aplicados orientados al apoyo al
diagnóstico dermatológico, con el correspondiente retorno previsto
sobre la práctica asistencial y la formación de profesionales. Su
uso en el presente trabajo se limita a experimentación parcial
sobre material clínico en español. El detalle exhaustivo del
pipeline de construcción, las distribuciones por L1/L2/L3, los
splits, el texto narrativo y los campos de validación experta se
desarrolla en el anexo~\ref{anexo:pipeline}.

\section{Ontología jerárquica L1/L2/L3}
\label{sec:ontologia}

La heterogeneidad de las taxonomías diagnósticas entre datasets ---
desde la categorización binaria melanoma/no-melanoma de Derm7pt y
HIBA hasta las 114 patologías incluidas en Fitzpatrick17k---
dificulta una comparación directa entre evaluaciones. Para permitir
comparaciones entre datasets y facilitar el desarrollo de un
clasificador unificado sobre un dominio común, el presente trabajo
emplea una ontología jerárquica de tres niveles diseñada por el autor
y revisada por la Dra.~R.~Taberner.

\subsection{Estructura jerárquica}
\label{sec:ontologia-estructura}

La ontología organiza los diagnósticos en tres niveles jerárquicos
encajados.

\textbf{Nivel L1 (etiológico)} agrupa los diagnósticos en cuatro
grandes familias etiológicas que reflejan el mecanismo causal
predominante: patología inflamatoria, patología tumoral, patología
infecciosa y genodermatosis. Estas cuatro categorías cubren el
espectro diagnóstico habitual esperado en la práctica dermatológica.

\textbf{Nivel L2 (subcategoría clínica)} desglosa cada categoría L1
en subcategorías clínicas más específicas. El vocabulario completo de
la ontología contiene 43 entradas L2 (p.~ej.~\emph{Eccemas y
dermatitis} dentro de Patología inflamatoria, \emph{Tumores
malignos} dentro de Patología tumoral o \emph{Infecciones por virus}
dentro de Patología infecciosa).

De estas 43 entradas, cinco no presentan diagnósticos L3 asociados:
\emph{Esclerosis tuberosa}, \emph{Neurofibromatosis tipo 1},
\emph{Incontinentia pigmenti}, \emph{Queratodermia palmoplantar
hereditaria} y \emph{Pseudoxantoma elástico}. Estas entidades se
consideran suficientemente específicas y mantienen su representación
directamente en el nivel L2, sin desagregación diagnóstica adicional.

\textbf{Nivel L3 (diagnóstico)} contiene los diagnósticos
individuales utilizados por el sistema de clasificación. El
vocabulario completo incorpora 367 entradas L3.

\subsection{Códigos externos asociados}
\label{sec:ontologia-codigos}

Cada entrada L3 de la ontología lleva asociados códigos de dos
nomenclaturas externas reconocidas: SNOMED CT
(\emph{Systematized Nomenclature of Medicine -- Clinical Terms}) y
CIE-10 (\emph{Clasificación Internacional de Enfermedades, décima
revisión}).

La incorporación de estos códigos persigue tres objetivos:
facilitar el mapeo bidireccional entre la ontología propia y las
taxonomías clínicas utilizadas habitualmente en el entorno
hospitalario, habilitar la interoperabilidad con sistemas de
información clínica y reducir el esfuerzo necesario para incorporar
nuevos diagnósticos mediante el respaldo de nomenclaturas externas
estables.

\subsection{Mapeo de los datasets a la ontología}
\label{sec:ontologia-mapeo}

Once de los doce datasets externos se mapean a la ontología
jerárquica mediante una correspondencia explícita entre cada
etiqueta nativa del dataset y la entrada L3 (o, en su defecto, L2)
correspondiente. Este mapeo permite construir un corpus armonizado
agregado de 72\,654 imágenes etiquetadas sobre vocabulario común,
que constituye la base del clasificador unificado descrito en el
capítulo~\ref{cap:caminos}. SkinCon, al estar anotado con conceptos
y no con diagnósticos categóricos, no participa de este mapeo a
nivel diagnóstico.

\section{Limitaciones del conjunto de datos del trabajo}
\label{sec:limitaciones-datos}

El conjunto de datasets externos descrito en este capítulo presenta
tres limitaciones que conviene declarar explícitamente, dado que
condicionan la interpretación de los resultados del
capítulo~\ref{cap:resultados}.

\textbf{(i)~Solapamiento con Derm1M.} Tres de los doce datasets
externos contrastados (Dermnet, HIBA, MSKCC) presentan solapamiento
documentado con el corpus de preentrenamiento Derm1M empleado por
DermLIP. Las evaluaciones sobre estos tres datasets se interpretan
con caveat explícito en el capítulo~\ref{cap:resultados}.

\textbf{(ii)~Sesgo geográfico del material.} Los doce datasets
externos provienen mayoritariamente de instituciones de Estados
Unidos, Australia, Europa central y Brasil. La población
hispanohablante europea no está representada de forma específica en
los conjuntos de evaluación.

\textbf{(iii)~Heterogeneidad de las taxonomías diagnósticas
nativas.} Los doce datasets emplean taxonomías propias incompatibles
entre sí (binarias en Derm7pt, HIBA y MSKCC; multiclase con número
variable en HAM10000, BCN20000, PAD-UFES-20, Dermnet y WSI patches;
114 patologías en Fitzpatrick17k). La ontología jerárquica
introducida en la sección~\ref{sec:ontologia} mitiga esta
limitación al armonizar el conjunto agregado, pero no la elimina:
cada dataset conserva su grano original y la comparación
inter-dataset requiere acumulación en un nivel L superior al de la
etiqueta nativa más fina.

El corpus propio DermapixelAI 1.0, al no ser objeto de evaluación
del trabajo, presenta limitaciones específicas (desbalance
estructural por modalidad, cola larga diagnóstica, cobertura
incompleta de la ontología y cobertura parcial por subcategoría en
los splits) que se documentan en detalle en el
anexo~\ref{anexo:pipeline}.

Estas limitaciones no invalidan los resultados del trabajo pero
delimitan el dominio de validez de las cifras reportadas en el
capítulo~\ref{cap:resultados} y las direcciones de continuación
recogidas en el capítulo~\ref{cap:caminos}.
